



\begin{landscape}

{\small                          % 字号调小，适配宽表
\setlength{\tabcolsep}{5pt}      % 列间距微调

\begin{longtable}{%
  L{2.0cm}   % Pathway
  L{2.8cm}   % Mechanism
  L{2.5cm}   % Construction Cost
  L{6.5cm}   % Information Transformation
  L{3.5cm}   % Incrementality
  L{2.8cm}   % Output Verifiability
  L{2.2cm}   % 代表工作 / 详节
}

%── 标题（如处于第 8 章，LaTeX 会自动编号为 Table 8.1）──────
\caption{Pathway production profiles across five dimensions.}%
\label{tab:pathway-profiles}\\

%── 首页页眉 ──────────────────────────────────────────────
\toprule
\textbf{Pathway} &
\textbf{Mechanism} &
\textbf{Construction Cost} &
\textbf{Information Transformation} &
\textbf{Incrementality} &
\textbf{Output Verifiability} &
\textbf{代表工作 / 详节} \\
\midrule
\endfirsthead

%── 续页页眉 ──────────────────────────────────────────────
\multicolumn{7}{c}{\tablename~\thetable\quad（续表）}\\[3pt]
\toprule
\textbf{Pathway} &
\textbf{Mechanism} &
\textbf{Construction Cost} &
\textbf{Information Transformation} &
\textbf{Incrementality} &
\textbf{Output Verifiability} &
\textbf{代表工作 / 详节} \\
\midrule
\endhead

%── 续页页脚 ──────────────────────────────────────────────
\midrule
\multicolumn{7}{r}{\small 接下页} \\
\endfoot

%── 末页页脚 ──────────────────────────────────────────────
\bottomrule
\endlastfoot

%── P1 ───────────────────────────────────────────────────
P1 Narrative Abstraction &
LLM inference &
低·可全自动 &
\textbf{保留}: 模式与洞察 /
\textbf{损失}: 步骤与时序(by-design) /
\textbf{注入}: 跨轨迹综合 insight(源自模型先验) &
增量(append); dynamic store 有 merge/prune 级联 &
直接阅读判断(主观) &
Reflexion [Shi23b]; \S3.1 \\
\midrule
%── P2 ───────────────────────────────────────────────────
P2 Schematic Formalization &
LLM inference + 可选 execution 验证 &
中·需可执行环境 &
\textbf{保留}: 过程结构 /
\textbf{损失}: NL nuance/隐含假设(by-design) /
\textbf{注入}: 显式依赖结构 &
增量(新 artifact 入库) &
execution test(客观, coverage-limited) &
Voyager [Wan23c]; \S3.2 \\
\midrule
%── P3 ───────────────────────────────────────────────────
P3 Latent-Space Transformation &
cache: gradient-free / trained: gradient &
cache 低·仅 session; trained 中 &
\textbf{保留}: 计算状态 /
\textbf{损失}: 符号可读性(by-design) /
\textbf{注入}: token 未显式编码的统计规律 &
cache 即时重建; trained 需 batch 重训 &
无直接验证·仅 ablation 间接 &
LAG [Che25b](cache)/
ReasonCACHE [Gup26](trained); \S4 \\
\midrule
%── P4 ───────────────────────────────────────────────────
P4 Evaluator Internalization &
gradient training &
高·需标注或明确成败信号 &
\textbf{保留}: 评估偏好 /
\textbf{损失}: 校准来源与单条贡献
(mix: by-design + by-limitation) &
batch 重训; 新标注需重训以保校准 &
无直接验证·hold-out calibration check &
Let's Verify Step by Step [Lig23]; \S5.1 \\
\midrule
%── P5 ───────────────────────────────────────────────────
P5 Policy Internalization &
gradient training (SFT/RL) &
高·需大规模高质量 trajectory &
\textbf{保留}: 决策模式 /
\textbf{损失}: 单条贡献不可区分(by-design)
+ 欠拟合/reward hacking 风险(by-limitation) &
batch 重训(可热启动); RL 可在线但漂移衰减旧经验 &
仅 behavioral (rollout 间接比较) &
SWE-Gym [Pan24]; \S5.2 \\
\midrule
%── P6 ───────────────────────────────────────────────────
P6 Evaluator-Driven Optimization &
gradient training (RLHF/DPO) &
高·叠加 P4 成本 &
\textbf{保留}: 偏好迁移入 policy /
\textbf{注入}: trajectory-only 学不到的细粒度偏好 /
\textbf{损失}: 经 P4+P6 双重压缩 &
online/iterative·漂移敏感 &
behavioral·循环验证风险
(Evaluator $\rightleftarrows$ Policy) &
Constitutional AI [Bai22b]; \S6.1 \\
\midrule
%── P7 ───────────────────────────────────────────────────
P7 Parametric Externalization &
从 policy 采样生成 &
中·依赖采样质量·须先有 trained policy &
\textbf{保留}: 权重中的能力经采样部分恢复为 token /
\textbf{损失}: 采样未覆盖部分 /
\textbf{注入}: hallucination 风险 &
可增量生成新样本·跟随源模型更新 &
产物可读但不可追溯权重成分·完备性不可验证 &
Explorer [Pah25]; \S6.2 \\

\end{longtable}

} % end \small

\end{landscape}